\subsection{One exclusion criterion for several geometric cases}
\label{sec:gradientexclusion}

The smooth case does not need a list of critical-point formulas.
The same argument also treats many singular leading forms. Its
hypotheses must hold in every complex direction. A bound only on
the original real integration space would not suffice.

We state the criterion for a polynomial exponential amplitude, then
specialize it to the Gaussian. Norms are Hermitian norms on complex
affine space; the Gaussian quadratic remains holomorphic.

\begin{theorem}[A gradient criterion]\label{thm:gradientcriterion}
Let $P,q\in\C[x_1,\ldots,x_n]$, with $P$ nonconstant and
$\deg q\leq J$, where $J\geq1$. Suppose that outside a compact set
\begin{equation}\label{eq:gradientconditions}
 \|\nabla P(x)\|\geq c\|x\|^{J-1},\qquad
 \frac{\|\operatorname{Hess}P(x)\|}{\|\nabla P(x)\|^2}\leq C
\end{equation}
for some $c>0$ and $C<\infty$. For every fixed polynomial $A$,
no nonzero flat period of
\[
 A(x)e^{tP(x)-q(x)}\,dx_1\cdots dx_n
\]
is a finite exponential-polynomial sum
$\sum_\lambda e^{\lambda t}p_\lambda(t)$ with polynomial
$p_\lambda$. In particular it cannot be a nonzero constant.
\end{theorem}
\begin{proof}
We first justify the full cycle space and the tails. The gradient is
nonzero outside a compact set, so the affine critical set is compact
and algebraic, hence finite. The gradient bound is stronger than the
condition excluding vanishing cycles at infinity. The finite-critical
thimble decomposition therefore applies, with the whole local Milnor
homology at each isolated critical point, not just its Morse
approximation~\cite[Part~I, Sections~1--4]{Pham}.
When several critical points have the same value, take disjoint
Milnor balls and apply excision to their complement. This gives the
direct sum of their full local homology before any values are grouped.

Here the Gaussian or polynomial amplitude can be kept on these fixed
$P$-cycles. Choose a generic argument $\theta$ of $t$. Outside the
local critical neighbourhoods, lift a descending phase ray by
\[
 P(x)=\lambda-e^{-i\theta}s,\qquad
 \frac{dx}{ds}=-e^{-i\theta}
                    \frac{\overline{\nabla P(x)}}{\|\nabla P(x)\|^2}.
\]
For large $s$ the lift is outside the compact set, because its phase
value is unbounded. The first condition in
\eqref{eq:gradientconditions} gives
$\|x\|^J\leq C_1(1+s)$. The real differential of this vector field
has norm at most a constant times the second ratio in
\eqref{eq:gradientconditions}, hence is bounded. Transported
$(n-1)$-cycle volumes grow at most as $\exp(C_2s)$; the radial speed
is bounded as well. Also
$|A(x)e^{-q(x)}|\leq C_3(1+s)^M\exp(C_4s)$ for some $M$.
Thus $\exp(-|t|s)$ dominates the complete tail, uniformly for large
$|t|$ in a smaller sector. The same estimate applies to every fixed
polynomial amplitude.

We must also check that these are a full basis for the actual phase,
rather than selected admissible cycles. Put
\[
 F_\sigma=P-\sigma q/t,\qquad 0\leq\sigma\leq1.
\]
Since $\|\nabla q\|\leq C_5\|\nabla P\|$ outside the same compact
set, all large $|t|$ give
\[
 \|\nabla F_\sigma\|\geq\tfrac12\|\nabla P\|.
\]
The ratio involving the Hessian stays uniformly bounded too. Indeed
$\|\operatorname{Hess}q\|=O(\|x\|^{J-2})$ for $J\geq2$, while
that Hessian is zero for $J=1$.
All critical points and their values therefore remain in fixed compact
sets. On the preimage of a phase half-plane beyond those values, the
parameter field
\[
 W_\sigma=\frac{q}{t}
       \frac{\overline{\nabla F_\sigma}}{\|\nabla F_\sigma\|^2}
\]
satisfies $\partial_\sigma F_\sigma+dF_\sigma(W_\sigma)=0$ and
$\|W_\sigma\|\leq C_6(1+\|x\|)/|t|$. Its flow is complete for
$0\leq\sigma\leq1$. Extend it by a cutoff away from that half-plane.
It gives an isotopy of the ambient relative pairs defining the full
rapid-decay homology. This explicitly prevents an infinity class
from appearing during the deformation.

The fixed $P$-thimbles have uniformly descending tails for every
$F_\sigma$, by the preceding estimates. Truncate them beyond a common
large phase value. Their boundaries stay in the descending half-plane
throughout the isotopy, so they define its flat relative classes.
They are a full basis at $\sigma=0$ and remain one at $\sigma=1$.
Finite critical points may merge, and their critical values may
coincide; the argument uses the full relative pairs and does not
require a Morse family. The same cycles are flat as $t$ varies in
the chosen sector.

These improper chains and their homotopies are admissible for the
period pairing. One can integrate first on truncated chains and then
pass to the limit: the tail bound makes every boundary term tend to
zero. Indeed the differential of $W_\sigma$ grows at most as
$C(1+\|x\|^J)/|t|$; along the controlled tails this is
$C'(1+s)/|t|$. The homotopy volumes have at most this exponential
growth, still dominated by the descending factor for all large $|t|$.
Thus Stokes and differentiation under the integral hold throughout;
the relative-pair calculation does not depend on a convergence claim
for the original real Gaussian integral.

It remains to determine the local asymptotics at a finite critical
point with value $\lambda$. For an analytic amplitude germ, the
vanishing-cycle density is a finite sum of terms
$s^{\alpha+k}(\log s)^j$, with analytic coefficient series and
$\alpha>-1$. Its Laplace asymptotics consequently have the form
\[
 e^{\lambda t}\sum_{\alpha,k,j}
       c_{\alpha,k,j}\,t^{-\alpha-k-1}(\log t)^j,
 \qquad k\geq0.
\]
This is the local analytic result of
Pham~\cite[Part~II, Section~1.1, (1.5)--(1.6) and the footnote on
printed page~26]{Pham}; it permits degenerate isolated critical
points. Pham credits Maisonobe and Rombaldi for the local solution
comparison used there. The amplitude $A e^{-q}$ is analytic. Every nonzero
term has a strictly negative power of $t$ and tends to zero after
removing $e^{\lambda t}$. The tails outside the local neighbourhood
are exponentially negligible by the uniform bound already proved.

The finite-critical decomposition accounts for the full module.
Group equal critical values before separating the distinct formal
exponential factors. Within each such factor the displayed local
expansions contain only decaying powers and logarithms. A nonzero
polynomial prefactor cannot occur. The zero function is allowed when
all these contributions vanish. This proves the theorem.
\end{proof}

\begin{corollary}[Original Gaussian markers]\label{cor:gradientmarkers}
For the standard Gaussian quadratic $q=\sum x_i^2/2$, suppose
\eqref{eq:gradientconditions} holds with $J=2$. Then every original
marked series $\Phi_A$ is either zero or is not a finite
exponential-polynomial sum. In particular $P$ is neither pure nor safe.
\end{corollary}
\begin{proof}
The exact top de Rham quotient and connection are
\[
 M=\C(t)[x]\Big/\sum_i(\partial_i+tP_i-x_i)\C(t)[x],
 \qquad\nabla[A]=[A_t+PA].
\]
Gaussian integration by parts makes $[A]\mapsto\Phi_A$ horizontal
on this vector-space quotient. The spatial and parameter derivatives
commute; no ordinary Jacobian-ideal quotient is substituted.
If some $\Phi_A$ is a
nonzero finite exponential-polynomial sum, a constant-coefficient
differential operator isolates one exponential and kills its
polynomial prefactor, giving $c e^{\beta t}$, $c\ne0$.
Because
\[
 \partial_t\Phi_A=\Phi_{PA},
\]
this is still an original fixed polynomial marker. Replace $P$ by
$P-\beta$ and divide the marker by $c$. Its formal seed value is one.
Lemma~\ref{lem:seedfunctional} gives an actual constant period for
that marker, contradicting Theorem~\ref{thm:gradientcriterion}.
Its comparison hypotheses hold with base $t$ and phase $tP-q$ on
$\mathbb A^1\times\mathbb A^n$, after removing the finite exceptional
base set. No convergence of the original Gaussian series is assumed.
Finally $\Phi_1(0)=1$, so it is a nonzero series; purity or safety
would make it a polynomial.
\end{proof}

\begin{corollary}[Smooth leading form, in every degree]
\label{cor:smoothleading}
If $\deg P=D\geq2$ and the homogeneous leading form satisfies
$\nabla F_D(v)=0$ only for $v=0$, then $P$ is not pure. In fact
Corollary~\ref{cor:gradientmarkers} holds for every original marker.
\end{corollary}
\begin{proof}
Compactness of the complex unit sphere gives
$\|\nabla P(x)\|\geq c\|x\|^{D-1}$ for large $\|x\|$.
Also $\|\operatorname{Hess}P\|=O(\|x\|^{D-2})$, so its ratio to
the squared gradient is $O(\|x\|^{-D})$. Both conditions for $J=2$
follow, with all lower terms retained.
\end{proof}

\begin{lemma}[The ordinary-node test]\label{lem:nodecoercivity}
Let $P=F_3+F_2+F_1+F_0$ in $n\geq2$ variables. Suppose its
projective singular set is finite and every singular point $[v]$
has $\ker\operatorname{Hess}F_3(v)=\C v$. If $F_2(v)\ne0$ at
every such point, then
$\|\nabla P(x)\|\geq c\|x\|$ outside a compact set. Consequently
$P$ is not pure under the Gaussian functional.
\end{lemma}
\begin{proof}
Otherwise choose an unbounded sequence with
$\|\nabla P(x)\|/\|x\|\longrightarrow0$. A projective limit $[v]$
is singular for $F_3$. Choose a fixed complement to $\C v$ and write
$x=sv+y$, where $y=o(s)$. With $H=\operatorname{Hess}F_3(v)$,
homogeneity gives the exact expansion
\[
 \nabla P(sv+y)=sHy+\nabla F_3(y)+s\nabla F_2(v)
                              +\nabla F_2(y)+\nabla F_1.
\]
The map $H$ is injective on the chosen complement. Divide by $s$.
The error from $\nabla F_3(y)/s$ is $o(\|y\|)$, so the equation
first forces $y$ to stay bounded. A limit $y_0$ satisfies
$Hy_0+\nabla F_2(v)=0$. Pair bilinearly with $v$. Symmetry and
$Hv=0$ kill the first term, whereas Euler's identity gives
$v\cdot\nabla F_2(v)=2F_2(v)$, a contradiction.
The Hessian of a cubic is $O(1+\|x\|)$, so the second condition
in \eqref{eq:gradientconditions} follows from the first. Apply
Corollary~\ref{cor:gradientmarkers}.
\end{proof}

\subsection{A common residual family for ordinary nodes}
\label{sec:ordinaryresidual}

For a pure cubic with only ordinary leading nodes, the preceding
lemma forces $F_2$ to vanish at some node. An isotropic such node
gives the frame already proved in Lemma~\ref{lem:frame}. Every
nonisotropic such node has the following single normal form.

\begin{lemma}\label{lem:ordinarynormal}
Let $[w]$ be an ordinary node of a ternary $F_3$, with $F_2(w)=0$
and $w^Tw\ne0$. An invertible affine change of coordinates puts
the pair consisting of $P$ and the Gaussian quadratic into
\begin{align*}
 P={}&T(uv+h)+a u^3+d v^3+b u^2+c v^2+e u+f v+k,\\
 q={}&\tfrac12(Au^2+Bv^2+mT^2+2Cuv+2DuT+2EvT)
                  +\ell_u u+\ell_v v+\ell_T T+q_0,
\end{align*}
where $m\ne0$ and $\det\operatorname{Hess}q\ne0$.
All Gaussian means and cross terms must be retained. Its leading
curve is irreducible nodal when $ad\ne0$, a transverse line and
smooth conic when exactly one of $a,d$ vanishes, and three
nonconcurrent lines when $a=d=0$.
\end{lemma}
\begin{proof}
Take $T$ in the direction $w$ and choose a complementary plane.
The singular-point equations make $F_3$ linear in $T$; its
coefficient is a nondegenerate binary quadratic because the node
is ordinary. Factor that quadratic and choose its two linear
factors as coordinates $u,v$. Thus $F_3=Tuv+B_3(u,v)$.
Replace $T$ by $T$ plus a linear form in $u,v$ to remove the
$u^2v$ and $uv^2$ terms of $B_3$.

There is no $T^2$ term in $F_2$, precisely because $F_2(w)=0$.
The full coefficient of $T$ is therefore
$uv+\alpha u+\beta v+h_0$. Translations of $u,v$ make it $uv+h$.
The cubic terms outside $T$ remain pure cubes. A constant shift
of $T$ removes the remaining quadratic $uv$ term. This gives the
displayed $P$. The direction of $T$ has not changed, so $m\ne0$;
invertibility of the affine change preserves the nondegenerate
Gaussian Hessian, while translations and the shifts of $T$ produce
the means and cross terms shown.

If $ad\ne0$, the polynomial $Tuv+a u^3+d v^3$ is irreducible:
an independent factor of a polynomial linear in $T$ would divide
both $uv$ and $a u^3+d v^3$, whose gcd is one. Its gradient shows
that its sole singular point is the stated ordinary node. If
$a=0,d\ne0$, it is $v(Tu+d v^2)$, with a smooth conic meeting
the line transversely at two points; the other one-zero case is
symmetric. If both vanish, it is the coordinate triangle $Tuv$.
\end{proof}

This consolidates three leading-curve rows into one residual family.
It is a reduction, not yet an exclusion of that family. In particular
the full quadratic $q$, including its means, cannot be replaced by
a convenient diagonal form during the remaining period argument.

\subsection{Two controls on the scope of the criterion}

Ordinary tameness of $P$ alone is insufficient for the Gaussian
conclusion. Let $P=\ell^Tx$, with $\ell=(1,i)/\sqrt2$ in two
variables. Its gradient is constant and nonzero, but
$\Phi_1=\exp(t^2\ell^T\ell/2)=1$. The phase $tP-q$ has its
critical point at $x=t\ell$, escaping with action zero. Thus an
isolated critical set for $P$, or the weaker gradient bound away
from zero, would lose exactly the contribution that purity needs.
Condition~\eqref{eq:gradientconditions} with $J=2$ prevents this.

The distinction between an original marker and a time-dependent
amplitude also remains essential. For
$P=(X_1^2+X_2^2)/2$ one has $\Phi_1=(1-t)^{-1}$ and
$(1-t)\Phi_1=1$. This phase satisfies the gradient criterion.
The unit in its rational differential module is represented by a
time-dependent amplitude. No original fixed polynomial marker has
a nonzero polynomial period, by Corollary~\ref{cor:gradientmarkers}.
The abstract module alone does not record that distinction.
